%-------------------------------------------------------------------------------
%	ABSTRACT PAGE
%-------------------------------------------------------------------------------

\addchaptertocentry{Хураангуй} % Хураангуйг гарчигт нэмэх

\begin{center}
{\scshape\Large \univname\par} % Их сургуулийн нэр
{\scshape\large \facname\par}\vspace{0.5cm} % Их сургуулийн нэр
{\huge\textbf{{Хураангуй}} \par}
\bigskip
{\Large{\ttitle} \par} % Тезисийн нэр
\bigskip

{\normalsize \shortname \par} % Зохиогчийн нэр
\addressname
\end{center}

\textit{\textbf{Түлхүүр үгс: \keywordnames}}
\bigskip


Энэхүү дипломын ажлын хүрээнд Монгол Улсын хууль эрх зүйн баримт бичгүүдэд суурилсан, хиймэл оюун ухааны RAG (Retrieval-Augmented Generation) архитектур бүхий ухаалаг асуулт–хариултын системийг боловсруулан хэрэгжүүлэв. Судалгааны ажлын гол зорилго нь хэрэглэгчийн асуултад хуулийн албан ёсны эх сурвалжид тулгуурласан, үнэн зөв, найдвартай хариулт өгөх механизмыг бүрдүүлэхэд оршино.

Системийн хэрэгжүүлэлтэд PDF болон текстэн өгөгдлийг автоматаар векторжуулах (Embedding), FAISS индекс ашиглан семантик хайлт хийх боломжийг бүрдүүлсэн бөгөөд Unsupervised learning-ийн clustering аргаар баримт бичгийг сэдэвчилсэн бүлэгт ангилах шийдлийг нэвтрүүлсэн. Мөн хэрэглэгчийн эргэх холбоонд тулгуурлан хайлтын чанарыг сайжруулах Supervised learning элементүүдийг нэгтгэж, системийн тасралтгүй суралцах чадварыг нэмэгдүүлэв. Гүйцэтгэлийн үнэлгээг Recall@k болон MRR хэмжүүрүүдээр тооцож, Reranker болон A/B тестээр оновчтой шийдлийг тодорхойлсон болно. Энэхүү ажил нь LLM технологийг бодит өгөгдлийн шинжилгээ, аналитик үнэлгээтэй хослуулсан бүрэн хэмжээний LegalTech шийдэл болсон юм.



